% ==============================================================================
% CAPÍTULO 3 - METODOLOGIA
% ==============================================================================

\chapter{Metodologia}
\label{cap:metodologia}

Para o desenvolvimento do projeto Oficina Master, foi adotada uma metodologia baseada em princípios ágeis, com entregas contínuas e feedback constante da equipe técnica da FK Oficina. O processo foi dividido em etapas de planejamento, levantamento de requisitos, desenvolvimento, testes e implantação.

A sistemática de trabalho envolveu o uso de ferramentas de versionamento e colaboração. O código-fonte foi gerenciado por meio do Git, utilizando a plataforma GitHub para revisões de código (*Pull Requests*) e integração contínua. Para a gestão das tarefas, utilizou-se o modelo *Kanban*, permitindo a visualização clara do fluxo de desenvolvimento (A fazer, Em desenvolvimento, Revisão e Concluído).

O ambiente de desenvolvimento local foi configurado utilizando ferramentas de conteinerização, garantindo a paridade entre os ambientes de desenvolvimento e produção. Para o backend, utilizou-se o Composer para gestão de pacotes PHP e o Artisan para automação de tarefas. No frontend, o NPM gerenciou as dependências do Vue.js, enquanto o Vite atuou como ferramenta de build ultra-rápida. A comunicação entre o frontend e o backend foi estabelecida via requisições assíncronas (AJAX) utilizando a biblioteca Axios, seguindo o padrão RESTful.
